\documentclass[11pt]{article}
\usepackage{amsmath}
\title{The travel cost of making a formation observation possible}
\begin{document}
\maketitle

A timely gap bit permits an advance that the same formation cannot
guarantee if it has to move to make the observation possible.
For a shortfall \(d\) in detection range, the exact loss in this
model is \(\sqrt{3d/2}\), not \(d\).

\paragraph{Source and motion law.}
Two single-file robots have the previous asymmetric-gap model
with \(a=V=1,A=4,W=3\), deadline 4, and a fixed physical gap band
\([1,8/5]\). Their initial gap is \(6/5\) (near) or \(7/5\) (far).
Both start at rest and must make one common advance \(D\) by
time 4, again at rest. In this study reversals are not admissible;
each velocity is nonnegative. The earlier optimum with a free
timely bit already has a forward-only realization, so that
comparison has not lost its witness.

The only gap evidence is one ideal range sample the controller
may schedule. It returns detection iff the instantaneous gap is
at most \(6/5-d\), with \(0<d\le3/800\). The reading is delivered
immediately, and a valid negative reading is distinct from
absence of a report. All commands before the sample must agree
in the two worlds. The measurement model deliberately omits
sensor noise, occlusions, and delivery latency. It is not a
claim about any real camera or communication system.

\paragraph{Exact upper bound.}
Write \(y=x_2-x_1\) for the difference of advances. Before the
sample every common history keeps \(y\in[-1/5,1/5]\).
At a distinguishing read \(q\), the far case cannot be
detected within that intersection band, and the near case
needs \(y(q)\le-d\). If there is no distinguishing read
before time 1, the near case must still satisfy the far case's
upper gap bound at 1. Hence \(x_2(1)\le1/2+1/5=7/10\),
which later gives \(D\le17/5\).

If the read is early, \(x_1(q)\le q^2/2\) and so
\(x_2(q)\le q^2/2-d\). Forward-only motion forces
\(q\ge\sqrt{2d}\). For \(q\le1/4\), a speed \(v\) at the
read requires \(x_2(q)\ge v^2/8\) by integrating the minimum
feasible velocity from rest. The remaining distance up to
time 1 is at most \(1-q-(1-v)^2/8\), attained by reaching
speed one as quickly as possible and cruising. Both position
and this bound are increasing in the allowed cut values.
Consequently
\[
 x_2(1)\le\frac78-q+\frac12\sqrt{q^2-2d}.
\]
The strongest value over those \(q\) is \(7/8-k\), where
\(k=\sqrt{3d/2}\). Indeed
\(q-\sqrt{q^2-2d}/2\ge k\) by squaring with \(q-k>0\);
the resulting inequality is \(3(q-4k/3)^2\ge0\).
The maximizing time \(4k/3\) lies in the early interval.
For \(q\in[1/4,1]\), the speed bound alone gives
\(x_2(1)\le q^2/2-d+1-q\le25/32-d<7/8-k\) in the declared
range of \(d\).

In the near case, robot 2 can add at most two units between
times 1 and 3. Robot 1's final acceleration bound gives
\(x_1(3)\ge D-1/2\), and near-case lower clearance gives
\(x_2(3)\ge x_1(3)-1/5\); so the final interval adds at most
\(7/10\) to robot 2's advance. Thus any policy has
\(D\le143/40-k\). This is above \(17/5\) for all the
specified \(d\); the late-read bound is not stronger.

\paragraph{An attaining active measurement.}
Let \(q=4k/3\). Robot 1 accelerates from time zero at one.
Robot 2 waits until \(3q/4\), then accelerates at four.
At the read both speeds are \(q\), with
\(x_1=q^2/2,\ x_2=q^2/8,\ y=-d\).
The controllers can now tell which physical gap they started with.

After a near reading keep robot 1's ramp until 1 and let robot 2
continue to accelerate until its speed cap. The leading gap at
1 is \(p=3/8-k\). In the far branch interpolate robot 2's
velocity after the read between \(v_1(t)=t\) and the near branch's
fast profile, with weight \((1/5+d)/(p+d)\) on the latter.
Relative displacement grows monotonically to \(1/5\). In both
branches from 1 to 3, robot 2 cruises at one and robot 1 adds
a triangular speed profile with area \(p+1/5\).
The final reversed ramps use lag \(1/5\) in the near branch
and \(p\) in the far branch. Both reach the asserted optimum.
All accelerations, speeds and gap extrema are verified in
the exact replay, including the initial negative excursion.

If the report is erased, the same prefix can preserve the
intersection band: use the far branch's initial ramp to \(1/5\),
add only \(2/5\) in the middle, and use lag \(1/5\) at the end.
This safely achieves \(17/5\), not the target with reliable
evidence. It counts as service only if that smaller advance
is independently allowed.

\paragraph{Matched comparisons.}
At \(d=3/800\) the read is at \(1/10\) and the maximum is
\(7/2=3.5\), compared with 3.575 when directly informed and
3.4 with no usable bit. A separate sampled shared-prefix LP
agrees at the named read and brackets nearby alternatives.
The continuous proof, not the sampled comparison, supplies the
global upper bound. This is a motion cost of acquiring an ideal
range reading under fixed laws, not a sensor characterization or
a win over an equally informed trajectory optimizer. No full
R3a, R6, or I1 acceptance follows.
\end{document}
